\documentclass[11pt,a4paper]{article}
\usepackage[T1]{fontenc}
\usepackage[utf8]{inputenc}
\usepackage{lmodern}
\usepackage[margin=27mm,headheight=15pt]{geometry}
\usepackage{microtype,amsmath,amssymb,amsthm,mathtools}
\usepackage{booktabs,longtable,array,xcolor,enumitem,listings,fancyhdr,titlesec,float}
\usepackage[breaklinks=true]{hyperref}
\usepackage[nameinlink,noabbrev]{cleveref}
\definecolor{ink}{HTML}{17324D}
\definecolor{accent}{HTML}{176B75}
\definecolor{soft}{HTML}{F2F5F7}
\hypersetup{colorlinks=true,linkcolor=ink,citecolor=accent,urlcolor=accent,
 pdftitle={Tetration, perfect powers, and frozen decimal digits},
 pdfauthor={ChatGPT},pdfsubject={A closed formula and elementary proof for an extremal tetration problem}}
\titleformat{\section}{\Large\bfseries\color{ink}}{\thesection}{0.65em}{}
\titleformat{\subsection}{\large\bfseries\color{ink}}{\thesubsection}{0.65em}{}
\pagestyle{fancy}
\fancyhf{}
\fancyhead[L]{\small\color{ink}Tetration and perfect-power minima}
\fancyhead[R]{\small Research article}
\fancyfoot[C]{\thepage}
\renewcommand{\headrulewidth}{0.3pt}
\setlength{\parskip}{4pt}
\setlength{\parindent}{1.2em}
\setlist{topsep=4pt,itemsep=2pt,parsep=0pt}
\newtheorem{theorem}{Theorem}[section]
\newtheorem{lemma}[theorem]{Lemma}
\newtheorem{proposition}[theorem]{Proposition}
\newtheorem{corollary}[theorem]{Corollary}
\theoremstyle{definition}
\newtheorem{definition}[theorem]{Definition}
\newtheorem{remark}[theorem]{Remark}
\newcommand{\N}{\mathbb N}
\newcommand{\Z}{\mathbb Z}
\newcommand{\vp}{\nu}
\newcommand{\V}{V}
\newcommand{\crt}{\operatorname{CRT}}
\newcommand{\lc}{\operatorname{lcm}}
\newcommand{\bits}{\operatorname{bits}}
\DeclareMathOperator{\ord}{ord}
\lstset{basicstyle=\small\ttfamily,backgroundcolor=\color{soft},
 frame=single,rulecolor=\color{soft},breaklines=true,columns=fullflexible,
 keepspaces=true,showstringspaces=false,tabsize=4,keywordstyle=\color{accent}\bfseries,
 commentstyle=\color{black!60},stringstyle=\color{ink}}

\crefname{theorem}{Theorem}{Theorems}
\crefname{lemma}{Lemma}{Lemmas}
\crefname{proposition}{Proposition}{Propositions}
\crefname{corollary}{Corollary}{Corollaries}
\crefname{section}{Section}{Sections}
\crefname{table}{Table}{Tables}

\begin{document}
\begin{titlepage}
\vspace*{9mm}
{\color{accent}\sffamily\bfseries RESEARCH ARTICLE\par}
\vspace{8mm}
{\Huge\bfseries\color{ink} Tetration, perfect powers,\\[3mm] and frozen decimal digits\par}
\vspace{8mm}
{\Large A closed formula for the smallest $n$th perfect power\\[2mm]
whose tetration has congruence speed $n$\par}
\vspace{9mm}
{\large 19 September 2026\par}
\vspace{8mm}
\hrule
\vspace{7mm}
\noindent\textbf{Abstract.}
Let $V(a)$ be the eventual number of additional decimal digits shared by
successive power towers with integer base $a>1$, where $10\nmid a$.
A 2025 MathOverflow question asks for the smallest $n$th perfect power
$\rho(n)=x^n>1$ with $V(x^n)=n$, and proposes that the minimizer always
ends in $5$ once $n\ge3$. We give a complete elementary proof of that
pattern and an explicit formula. Set $s=n-\nu_2(n)$. Then
$\rho(1)=2$, $\rho(2)=49$, and, for $n\ge3$,
\[
\rho(n)=B_s^n,\qquad
B_s=\begin{cases}
2^s-1,&s\equiv0\pmod4,\\
3\cdot2^s-1,&s\equiv1\pmod4,\\
2^s+1,&s\equiv2\pmod4,\\
3\cdot2^s+1,&s\equiv3\pmod4.
\end{cases}
\]
The key is a deterministic size barrier: every competing root not divisible
by $5$ must satisfy $x^2+1\ge5^{n-\nu_5(n)}$. This excludes all competitors
for $n\ge25$; five small Hensel-lifting certificates settle the remaining
exceptional comparisons. We also solve the stricter maximal-degree version:
only $n=3$ changes, from $25^3$ to $55^3$. Further results include the exact
choice between the two conjectured decimal tails, a two-parameter
extremal theorem, density formulas, and asymptotics. The archive contains
standard-library Python code, exact certificates, and reproducible tests.

\vspace{5mm}
\noindent\textbf{Scope and status.}
The argument below is a complete proof for the precisely stated arithmetic
problem, not a numerical extrapolation. The foundational congruence-speed
formula is prior work and is rederived for completeness. The consulted
question leaves the all-$n$ extremal pattern conjectural; an exhaustive
claim of bibliographic priority is not made. This report has not undergone
external peer review or proof-assistant verification.

\vfill
\noindent\small Prepared by ChatGPT in response to a mathematical research request.\\
Source-status check: 19 September 2026.
\end{titlepage}

{\setlength{\parskip}{0pt}\tableofcontents}
\clearpage

\section{The problem, its provenance, and the answer}
\label{sec:problem}

\subsection{Why this particular question?}
An integer power tower can be far too large to write down while its last
few digits are easy to study. This makes modular tetration a useful place
to seek a tractable problem: one can replace the enormous integers by
small local invariants, then turn a computational pattern into a proof.
The relevant objects here are ordinary integer tetrations
$a\uparrow\uparrow h$, not analytic interpolations in the height.

The eventual decimal congruence-speed formula was established in work of
Rip\`a \cite{ripa2020,ripa2021}; stable-digit counts were subsequently
studied by Rip\`a and Onnis \cite{ripaonnis2022}. Those results are not the
open problem attacked here. The later extremal question asks that the
base itself be an $n$th perfect power and that its speed equal the same
integer $n$.

Specifically, Rip\`a's February 2025 MathOverflow question asks for a closed
formula for this minimum. Alekseyev's answer gives 50 root values and
observes the apparent predominance of the final-digit-$5$ case. The
question also conjectures two fixed ten-digit tails for all $n>10$
\cite{mo2025}. A related paper, published in 2024 and deposited on arXiv in January
2026, constructs perfect powers with prescribed speeds. Existence is
different from the global minimality proved here \cite{ripa2026}.

The choice is deliberately narrow. No conclusion below is a proposed
resolution of analytic tetration, arbitrary hyperoperation interpolation,
or ordinal fast-growing-hierarchy questions. It is a concrete extremal
problem inside the arithmetic of tetration.

\subsection{Definitions and conventions}
For $a\ge2$, put
\begin{equation}
T_0(a)=1,\qquad T_{h+1}(a)=a^{T_h(a)}\quad(h\ge0).
\label{eq:tower}
\end{equation}
Thus $T_h(a)=a\uparrow\uparrow h$. For a nonzero integer $u$ and a prime
$p$, $\nu_p(u)$ is the largest $j\ge0$ such that $p^j\mid u$.
Define
\begin{equation}
\Delta_h(a)=T_{h+1}(a)-T_h(a),\qquad
D_h(a)=\min\{\nu_2(\Delta_h(a)),\nu_5(\Delta_h(a))\}.
\label{eq:delta}
\end{equation}
Exactly $D_h(a)$ trailing decimal positions agree between these two
towers. For $10\nmid a$, we shall prove that
\[
D_h(a)-D_{h-1}(a)
\]
is eventually a positive constant; this constant is $V(a)$.
For convenience only, set $V(1)=0$. Multiples of $10$ are excluded from
the extremal problem because their speed is not eventually a finite
constant.

\begin{definition}
For $n\ge1$, define
\begin{align}
r(n)&=\min\{x\in\N:x>1,\ 10\nmid x,\ V(x^n)=n\},\label{eq:rdef}\\
\rho(n)&=r(n)^n.\label{eq:rhodef}
\end{align}
The minimum will be shown to exist. Since $x\mapsto x^n$ is strictly
increasing on the positive integers, minimizing $x$ and minimizing $x^n$
are equivalent.
\end{definition}

An $n$th perfect power need not have \emph{maximal} perfect-power degree
$n$. For example, $25^3=5^6$ is an admissible cube under
\eqref{eq:rhodef}. We solve the maximal-degree variant separately in
\cref{sec:strict} rather than silently changing the question.

\subsection{Main theorem}
\begin{theorem}[Closed formula]
\label{thm:main}
We have $r(1)=2$ and $r(2)=7$. For every $n\ge3$, let
\[
s=n-\nu_2(n)\ge2.
\]
Then
\begin{equation}
r(n)=B_s,
\qquad
B_s=\begin{cases}
2^s-1,&s\equiv0\pmod4,\\
3\cdot2^s-1,&s\equiv1\pmod4,\\
2^s+1,&s\equiv2\pmod4,\\
3\cdot2^s+1,&s\equiv3\pmod4.
\end{cases}
\label{eq:main}
\end{equation}
Consequently $r(n)\equiv\rho(n)\equiv5\pmod{10}$ for all $n\ge3$.
\end{theorem}

The first root values are
\[
2,\ 7,\ 25,\ 5,\ 95,\ 95,\ 385,\ 95,\ 1535,\ 1535,\ 6145,\ 1025,\ldots.
\]
The proof has two genuinely different parts. Finding the smallest
admissible multiple of $5$ is an elementary congruence calculation.
Proving that it beats \emph{every} other residue class requires the size
barrier in \cref{sec:barrier}. Checking many initial terms would not
replace that second part.

\section{From power towers to two local valuations}
\label{sec:local}

We give the details so that the main theorem does not rely on treating a
large published case distinction as a black box. The resulting compact
formula is equivalent to the known decimal congruence-speed formula
\cite{ripa2021,ripaonnis2022}.

\subsection{The elementary lifting facts}
\begin{lemma}[Lifting near $1$]
\label[lemma]{lem:lte}
Let $u>1$ and $m\ge1$ be integers. If $u\equiv1\pmod5$, then
\begin{equation}
\nu_5(u^m-1)=\nu_5(u-1)+\nu_5(m).
\label{eq:lte5}
\end{equation}
If $u\equiv1\pmod4$, then
\begin{equation}
\nu_2(u^m-1)=\nu_2(u-1)+\nu_2(m).
\label{eq:lte2}
\end{equation}
In addition, if $u$ is odd and $m$ is even, then
\begin{equation}
\nu_2(u^m-1)=\nu_2(u-1)+\nu_2(u+1)+\nu_2(m)-1.
\label{eq:lte2general}
\end{equation}
\end{lemma}
\begin{proof}
For either prime $p$, if $p\nmid q$ and $u\equiv1\pmod p$, the factorization
\[
\frac{u^q-1}{u-1}=1+u+\cdots+u^{q-1}\equiv q\not\equiv0\pmod p
\]
shows that raising to the $q$th power leaves the valuation unchanged.
It therefore suffices to understand raising to the $p$th power.

For $p=5$, write $u=1+5^j w$ with $5\nmid w$ and $j\ge1$. In the binomial
expansion of $u^5-1$, the first term has valuation $j+1$, whereas each
other term has larger valuation. Thus
$\nu_5(u^5-1)=j+1$. Repeating this step proves \eqref{eq:lte5}.
For $p=2$ with $u\equiv1\pmod4$, the factorization
$u^2-1=(u-1)(u+1)$ gives $\nu_2(u+1)=1$, so each squaring increases the
valuation by exactly one. This proves \eqref{eq:lte2}.
Finally, apply \eqref{eq:lte2} to $u^2\equiv1\pmod8$ and exponent $m/2$:
\[
\nu_2(u^m-1)=\nu_2(u^2-1)+\nu_2(m/2),
\]
which is \eqref{eq:lte2general}.
\end{proof}

For odd $a>1$ and, respectively, $5\nmid a>1$, define
\begin{equation}
c_2(a)=\nu_2(a^2-1)-1,
\qquad c_5(a)=\nu_5(a^4-1).
\label{eq:cs}
\end{equation}
For odd $a$, one of $a-1,a+1$ has valuation exactly one at $2$, so
\begin{equation}
c_2(a)=\max\{\nu_2(a-1),\nu_2(a+1)\}\ge2.
\label{eq:c2max}
\end{equation}
The functions $c_2,c_5$ are only used on the domains just specified.
In particular, no valuation of zero is needed.

\subsection{The recurrence for tower differences}
For $h\ge1$, direct factorization of \eqref{eq:tower} gives
\begin{equation}
\Delta_h=a^{T_{h-1}}\bigl(a^{\Delta_{h-1}}-1\bigr).
\label{eq:deltarec}
\end{equation}
If a prime $p$ divides $a$, the second factor is a $p$-adic unit. Hence
\begin{equation}
\nu_p(\Delta_h)=\nu_p(a)T_{h-1}
\quad\text{when }p\mid a.
\label{eq:nonunit}
\end{equation}
This term eventually dominates any affine function of $h$.
Indeed, $T_h\ge2^h$: for an integer $u\ge1$, $2^u\ge2u$, and induction
in \eqref{eq:tower} gives $T_{h+1}\ge2T_h$.

If $a$ is odd, all $T_h$ are odd, so every $\Delta_h$ is even.
By \cref{lem:lte}, \eqref{eq:deltarec} then implies
\begin{equation}
\nu_2(\Delta_h)=\nu_2(\Delta_{h-1})+c_2(a)
=\nu_2(a-1)+h c_2(a).
\label{eq:2affine}
\end{equation}
This identity holds for every $h\ge0$.

Suppose instead that $5\nmid a$. Once $h\ge3$, we have
$4\mid\Delta_{h-1}$. To see this when $a$ is odd, use
\eqref{eq:2affine} and $c_2(a)\ge2$. When $a$ is even, $T_j$ is divisible
by $4$ for every $j\ge2$, and so are the corresponding differences.
Now apply \eqref{eq:lte5} to $u=a^4$ and $m=\Delta_{h-1}/4$:
\begin{equation}
\nu_5(\Delta_h)=\nu_5(\Delta_{h-1})+c_5(a)
\qquad(h\ge3).
\label{eq:5affine}
\end{equation}
Thus the $5$-adic valuations are eventually affine with slope $c_5(a)$.

\begin{theorem}[Compact congruence-speed formula]
\label{thm:speed}
For $a>1$ and $10\nmid a$, the constant $V(a)$ exists and satisfies
\begin{equation}
V(a)=\begin{cases}
c_5(a),&2\mid a,\\
c_2(a),&5\mid a,\\
\min\{c_2(a),c_5(a)\},&\gcd(a,10)=1.
\end{cases}
\label{eq:speed}
\end{equation}
\end{theorem}
\begin{proof}
When $a$ is even, the $2$-adic term in $D_h$ has the rapidly growing form
\eqref{eq:nonunit}, while the $5$-adic term is eventually affine with
slope $c_5(a)$. Their minimum therefore eventually equals the latter.
The case $5\mid a$ is symmetric, with slope $c_2(a)$.
When $a$ is coprime to $10$, both terms are eventually affine. The minimum
of two affine functions is eventually affine with slope the smaller of
the two slopes; if the slopes are equal, only the intercepts matter.
This proves \eqref{eq:speed} and existence of $V(a)$.
\end{proof}

\begin{remark}[Transient behavior is not ignored]
The theorem concerns eventual slopes, not every height. For example,
for $a=5$ the first values $D_0,D_1,D_2,D_3,D_4$ are
$0,1,5,8,10$; only later is each increment $2$. For $a=807$, the
minimum switches between local affine contributions before settling at
slope $3$. The verification code checks tower residues directly, rather
than assuming that a low-height observed difference must equal $V(a)$.
\end{remark}

\section{Perfect powers shift the local speeds}
\label{sec:powers}

\begin{proposition}[Power-transfer law]
\label[proposition]{prop:power}
For $x>1$, $10\nmid x$, and $k\ge1$,
\begin{equation}
V(x^k)=\begin{cases}
c_5(x)+\nu_5(k),&2\mid x,\\
c_2(x)+\nu_2(k),&5\mid x,\\
\min\{c_2(x)+\nu_2(k),c_5(x)+\nu_5(k)\},&\gcd(x,10)=1.
\end{cases}
\label{eq:powertransfer}
\end{equation}
\end{proposition}
\begin{proof}
For odd $x$, apply \eqref{eq:lte2} to $u=x^2$ to obtain
\[
c_2(x^k)=\nu_2(x^{2k}-1)-1
=\nu_2(x^2-1)-1+\nu_2(k)=c_2(x)+\nu_2(k).
\]
For $5\nmid x$, apply \eqref{eq:lte5} to $u=x^4$:
\[
c_5(x^k)=\nu_5(x^{4k}-1)=c_5(x)+\nu_5(k).
\]
Substitute into \cref{thm:speed}.
\end{proof}

This identity is the main computational simplification. Evaluating
$V(x^k)$ requires no $k$th power, much less any tower with base $x^k$.
Only the valuations of $x\pm1$, $x^2+1$, and $k$ enter.

For the diagonal problem $k=n$, write
\begin{equation}
s=n-\nu_2(n),\qquad t=n-\nu_5(n).
\label{eq:st}
\end{equation}
If $5\mid x$ and $x$ is odd, then
\begin{equation}
V(x^n)=n\quad\Longleftrightarrow\quad c_2(x)=s.
\label{eq:fivecondition}
\end{equation}
On the other hand, every admissible root with $5\nmid x$ necessarily
satisfies
\begin{equation}
c_5(x)\ge t.
\label{eq:necessaryfive}
\end{equation}
The second assertion holds whether $x$ is even or odd. For odd $x$ it is
only a necessary condition, which is enough for a lower bound.

For $n\ge3$, $s\ge2$. This follows immediately for odd $n$; for even
$n=2^j u$ with $u$ odd, $n-j\ge2$, including the smallest case $n=4$.
Also $t\ge1$ for all positive $n$.

\section{The smallest root divisible by five}
\label{sec:restricted}

\begin{lemma}[Exact $2$-adic congruence classes]
\label[lemma]{lem:c2class}
For an odd $x>1$ and $s\ge2$,
\begin{equation}
c_2(x)=s\quad\Longleftrightarrow\quad
x\equiv2^s-1\ \text{or}\ 2^s+1\pmod{2^{s+1}}.
\label{eq:exactclasses}
\end{equation}
\end{lemma}
\begin{proof}
By \eqref{eq:c2max}, the larger of the valuations of $x-1,x+1$ must be
exactly $s$. This is equivalent to one of those two integers being
$2^s$ times an odd number. The other has valuation one. This yields
exactly the two stated classes.
\end{proof}

\begin{proposition}[Restricted minimum]
\label[proposition]{prop:restricted}
For $s\ge2$, the least odd positive multiple of $5$ satisfying
$c_2(x)=s$ is $B_s$ from \eqref{eq:main}. In particular,
\begin{equation}
5\mid B_s,\qquad c_2(B_s)=s,\qquad
2^s-1\le B_s\le3\cdot2^s+1.
\label{eq:Bbounds}
\end{equation}
\end{proposition}
\begin{proof}
By \cref{lem:c2class}, every candidate has the form
$x=j2^s+\varepsilon$, where $j\ge1$ is odd and
$\varepsilon\in\{-1,1\}$. The powers of $2$ modulo $5$ have period four.
The smallest admissible choices are therefore
\begin{center}
\begin{tabular}{ccccc}
\toprule
$s\bmod4$&$2^s\bmod5$&smallest odd $j$&$\varepsilon$&$B_s$\\
\midrule
0&1&1&$-1$&$2^s-1$\\
1&2&3&$-1$&$3\cdot2^s-1$\\
2&4&1&$+1$&$2^s+1$\\
3&3&3&$+1$&$3\cdot2^s+1$\\
\bottomrule
\end{tabular}
\end{center}
In the second and fourth rows, neither sign works for $j=1$, and $j=2$
is forbidden because the valuation must be exact. In all rows the
listed $j$ works; every larger eligible odd $j$ produces a larger $x$.
The valuation and bounds follow directly.
\end{proof}

For every $n\ge3$, this constructs an admissible $n$th power
$B_s^n$ and proves that no smaller multiple-of-$5$ root can work.
There remains a global question: could a root in a different decimal
class be even smaller?

\section{A deterministic barrier against all other roots}
\label{sec:barrier}

\begin{lemma}[Quadratic size barrier]
\label[lemma]{lem:barrier}
Let $x>1$, $5\nmid x$, and $t\ge1$. If $c_5(x)\ge t$, then
\begin{equation}
x^2+1\ge5^t.
\label{eq:barrier}
\end{equation}
\end{lemma}
\begin{proof}
Factor
\[
x^4-1=(x-1)(x+1)(x^2+1).
\]
Exactly one of the three factors is divisible by $5$. Indeed,
$x\equiv1,-1,2,-2\pmod5$ determines which factor it is, and two of these
factors cannot simultaneously be divisible by $5$. Thus all the factor
$5^t$ divides one of them. Each is positive and at most $x^2+1$ when
$x\ge2$, proving the result.
\end{proof}

The factorization is decisive. Merely using $5^t\mid x^4-1$ would give
only a fourth-root lower bound, too weak to beat a candidate of order
$2^n$. Isolating the factor $x^2+1$ raises the lower bound to approximately
$5^{n/2}$, which eventually dominates $2^n$.

\begin{proposition}[A sufficient global-minimality condition]
\label[proposition]{prop:criterion}
For $n\ge3$, let $s,t$ be as in \eqref{eq:st}. If
\begin{equation}
5^t>B_s^2+1,
\label{eq:criterion}
\end{equation}
then $r(n)=B_s$.
\end{proposition}
\begin{proof}
The root $B_s$ is admissible and is least among roots divisible by $5$.
A competing root $x\le B_s$ not divisible by $5$ would have $c_5(x)\ge t$
by \eqref{eq:necessaryfive}. But \cref{lem:barrier} would imply
$5^t\le x^2+1\le B_s^2+1$, a contradiction.
\end{proof}

\begin{lemma}[Uniform cutoff]
\label[lemma]{lem:cutoff}
Condition \eqref{eq:criterion} holds for every $n\ge25$.
\end{lemma}
\begin{proof}
Since $5^{\nu_5(n)}\le n$ and $s\le n$,
\[
5^t\ge\frac{5^n}{n},\qquad
B_s^2+1\le(3\cdot2^n+1)^2+1.
\]
Put $H_n=(3\cdot2^n+1)^2+1$. Direct expansion gives
\[
H_{n+1}=36\cdot4^n+12\cdot2^n+2
<4H_n.
\]
Consequently, for $n\ge5$,
\[
\frac{5^{n+1}/((n+1)H_{n+1})}{5^n/(nH_n)}
>\frac{5n}{4(n+1)}>1.
\]
The starting inequality at $n=25$ is the exact integer calculation
\begin{align*}
5^{25}-25H_{25}
&=298023223876953125-253327484072755250\\
&=44695739804197875>0.
\end{align*}
Thus $5^n/n>H_n$ for every $n\ge25$, proving the claim.
\end{proof}

Notice that no assumption about random $5$-adic digits occurs here.
The lower bound remains valid even if the residues of square roots of
$-1$ have extremely atypical digit patterns.

\section{The finite range: exact, inspectable certificates}
\label{sec:finite}

\subsection{The routine comparisons}
The uniform cutoff leaves only $3\le n\le24$. In
\cref{tab:small}, $Q_n=\lfloor5^t/(B_s^2+1)\rfloor$. A positive $Q_n$
certifies \eqref{eq:criterion}: equality $5^t=B_s^2+1$ cannot occur,
because $5\mid B_s$.

\begin{table}[H]
\centering
\small
\begin{tabular}{rrrrr}
\toprule
$n$&$s$&$t$&$B_s$&$Q_n$\\
\midrule
3&3&3&25&0\\
4&2&4&5&24\\
5&5&4&95&0\\
6&5&6&95&1\\
7&7&7&385&0\\
8&5&8&95&43\\
9&9&9&1535&0\\
10&9&9&1535&0\\
11&11&11&6145&1\\
12&10&12&1025&232\\
13&13&13&24575&2\\
14&13&14&24575&10\\
15&15&14&98305&0\\
16&12&16&4095&9099\\
17&17&17&393215&4\\
18&17&18&393215&24\\
19&19&19&1572865&7\\
20&18&19&262145&277\\
21&21&21&6291455&12\\
22&21&22&6291455&60\\
23&23&23&25165825&18\\
24&21&24&6291455&1505\\
\bottomrule
\end{tabular}
\caption{The finite comparison range. Only $n=3,5,7,9,10,15$ require
more than the quadratic size barrier. All entries are exact integers.}
\label{tab:small}
\end{table}

\subsection{A two-root lifting lemma}
\begin{lemma}
\label[lemma]{lem:hensel}
For every $k\ge1$, $z^2\equiv-1\pmod{5^k}$ has exactly two roots. Starting
with either root modulo $5$, each root has exactly one lift to a root
modulo $5^{k+1}$. The two roots are negatives of each other.
\end{lemma}
\begin{proof}
Modulo $5$, the roots are $2$ and $3$. Suppose $r^2+1=5^k q$.
Every lift is $r+j5^k$, with $j\in\{0,1,2,3,4\}$. Modulo $5^{k+1}$,
\[
(r+j5^k)^2+1\equiv5^k(q+2rj).
\]
Since $5\nmid2r$, there is a unique $j$ modulo $5$ satisfying
$q+2rj\equiv0\pmod5$. Every root at the next level reduces to a root
at the previous level, so no others appear. Negation interchanges the
two initial roots and preserves the equation at every level.
\end{proof}

For a root $x>1$ with $c_5(x)\ge t\ge k$, the factorization in
\cref{lem:barrier} gives one of
\begin{equation}
x\equiv1,-1\pmod{5^k},\qquad x^2\equiv-1\pmod{5^k}.
\label{eq:fourclasses}
\end{equation}
The smallest permitted $x>1$ in the first pair of classes is $5^k-1$.
If $d_k$ is the smaller root of $-1$ modulo $5^k$, the second pair
requires $x\ge d_k$. Thus
\begin{equation}
x\ge\min\{5^k-1,d_k\}.
\label{eq:exactlower}
\end{equation}

\subsection{Five certificates settle six indices}
The following table is enough; it does not search through a large
interval of candidate roots.
\begin{table}[H]
\centering
\small
\begin{tabular}{crrrrr}
\toprule
$n$&$B_s$&$k$&$5^k$&$d_k$&$(d_k^2+1)/5^k$\\
\midrule
3&25&3&125&57&26\\
5&95&4&625&182&53\\
7&385&5&3125&1068&365\\
9,10&1535&7&78125&32318&13369\\
15&98305&8&390625&110443&31226\\
\bottomrule
\end{tabular}
\caption{Certificates for the exceptional comparisons. In each row,
$k\le t$, $B_s<d_k<5^k/2$, and $B_s<5^k-1$.}
\label{tab:certs}
\end{table}
Each quotient in the last column is integral, so $d_k$ is a root of
$-1$ modulo $5^k$. It is less than half the modulus, and
\cref{lem:hensel} says the only other root is $5^k-d_k$. Therefore it
really is the smaller root. The inequalities in the caption, combined
with \eqref{eq:exactlower}, exclude every root $x\le B_s$ not divisible
by $5$ in each exceptional case.

\begin{proof}[Completion of \cref{thm:main}]
For $n\ge25$, use \cref{lem:cutoff,prop:criterion}. For $3\le n\le24$,
use \cref{tab:small} when $Q_n>0$ and \cref{tab:certs} otherwise. The
restricted minimum is always supplied by \cref{prop:restricted}.
For $n=1$, $V(2)=\nu_5(15)=1$, so the smallest possible root is $2$.
For $n=2$, direct application of \eqref{eq:powertransfer} gives
\[
\begin{array}{c|rrrrrr}
x&2&3&4&5&6&7\\ \hline
V(x^2)&1&1&1&3&1&2.
\end{array}
\]
Hence $r(2)=7$. Finally $B_s$ is odd and divisible by $5$, and its positive
powers all end in $5$.
\end{proof}

\section{The stricter maximal-degree problem}
\label{sec:strict}

The original discussion also asks about powers of \emph{exact} degree
$n$, meaning that the answer cannot also be a $d$th power with $d>n$
\cite{mo2025}. To avoid ambiguity, define the maximal perfect-power
degree of an integer $a>1$ to be
\[
\operatorname{ppdeg}(a)=\max\{d\ge1:a=y^d\text{ for an integer }y>1\}.
\]
If $a=\prod p^{e_p}$, this maximum is $\gcd_p e_p$. Consequently,
\begin{equation}
\operatorname{ppdeg}(x^n)=n
\quad\Longleftrightarrow\quad x\text{ is not a perfect power of degree }>1.
\label{eq:ppdeg}
\end{equation}
Let $\rho_{\mathrm{exact}}(n)$ be the minimum subject to
$V(a)=n$ and $\operatorname{ppdeg}(a)=n$, with $10\nmid a$.

\begin{lemma}[The restricted roots are almost never perfect powers]
\label[lemma]{lem:primitive}
For $s\ge2$, $B_s$ is a perfect power of degree greater than one if and
only if $s=3$. The exceptional value is $B_3=25$.
\end{lemma}
\begin{proof}
Any perfect power is either a square or a $q$th power for an odd prime
$q$. Its root is odd because $B_s$ is odd.

First consider $B_s=2^s-1$, where $s\equiv0\pmod4$. If $B_s=y^2$, then
$y^2+1=2^s$, impossible modulo $8$ for $s\ge2$. If $B_s=y^q$ with odd
$q\ge3$, factor $y^q+1$ as $(y+1)$ times
\[
y^{q-1}-y^{q-2}+\cdots-y+1.
\]
This second factor is odd and greater than one, so their product cannot
be a power of $2$.

Next consider $B_s=2^s+1$, where $s\equiv2\pmod4$. For an odd exponent,
$(y^q-1)/(y-1)$ is odd and greater than one, again contradicting a power
of $2$. For a square, $(y-1)(y+1)=2^s$ consists of two powers of $2$
differing by $2$. They must be $2$ and $4$, giving $s=3$, outside this
case.

Suppose $B_s=3\cdot2^s-1$, where $s\equiv1\pmod4$ and $s\ge5$. A square
would give $y^2+1=3\cdot2^s$, impossible modulo $8$. For an odd prime
exponent, the odd factor
\[
y^{q-1}-y^{q-2}+\cdots-y+1
\]
would divide $3$. But it is at least $y^2-y+1\ge7$ for odd $y\ge3$.
This is impossible.

Finally suppose $B_s=3\cdot2^s+1$, where $s\equiv3\pmod4$. For an odd
prime exponent, the factor $(y^q-1)/(y-1)$ is odd, divides $3$, and is
at least $y^2+y+1\ge13$, a contradiction. For a square, put
$u=(y-1)/2$ and $v=(y+1)/2$. Then
\[
v-u=1,\qquad \gcd(u,v)=1,\qquad uv=3\cdot2^{s-2}.
\]
The odd member of the consecutive pair is either $1$ or $3$.
The pair $1,2$ has product not divisible by $3$. The remaining pairs are
$2,3$ and $3,4$, giving $s=3$ and $s=4$, respectively. The congruence
$s\equiv3\pmod4$ leaves only $s=3$, with $B_s=25$.
\end{proof}

\begin{theorem}[Exact-degree minimum]
\label{thm:strict}
For every $n\ge1$,
\begin{equation}
\rho_{\mathrm{exact}}(n)=
\begin{cases}
55^3=166375,&n=3,\\
\rho(n),&n\ne3.
\end{cases}
\label{eq:strict}
\end{equation}
\end{theorem}
\begin{proof}
The roots $r(1)=2,r(2)=7$ are not perfect powers. For $n\ge3$,
\cref{lem:primitive} says that only $s=n-\nu_2(n)=3$ can cause trouble.
Writing $j=\nu_2(n)$, this equation gives $n=j+3$. For $j\ge3$,
$2^j>j+3$, impossible because $2^j\mid n$. The cases $j=0,1,2$ leave
only $n=3$. Thus for $n\ne3$, the unrestricted minimizer already has
maximal degree $n$.

For $n=3$, $c_2(55)=3$, so $V(55^3)=3$. The root $55=5\cdot11$ is not
a perfect power. Among smaller odd multiples of $5$, the only one with
$c_2(x)=3$ is $25$, as can be checked from the list
$5,15,25,35,45$ or from \eqref{eq:exactclasses}. It is a square.
Every root not divisible by $5$ with $c_5(x)\ge3$ is at least $57$ by
\eqref{eq:exactlower} at $k=3$. Therefore no root below $55$ satisfies
both required conditions.
\end{proof}

This strengthening needs no appeal to a deep theorem about consecutive
perfect powers. The four elementary factorizations above are sufficient.

\section{The decimal-tail conjecture, with the sign determined}
\label{sec:tails}

For $d\ge1$, let $E_d$ be the unique residue in $[0,10^d)$ satisfying
\begin{equation}
E_d\equiv0\pmod{5^d},\qquad E_d\equiv1\pmod{2^d}.
\label{eq:Ed}
\end{equation}
The Chinese remainder theorem gives existence and uniqueness, as well as
the explicit finite formula
\begin{equation}
E_d=5^d\bigl((5^d)^{-1}\bmod2^d\bigr).
\label{eq:Edformula}
\end{equation}
Here the inverse is represented in $\{0,\ldots,2^d-1\}$.
Since $E_d^2\equiv E_d$ modulo both prime powers, $E_d$ is idempotent
modulo $10^d$. The residues are compatible as $d$ grows.

\begin{theorem}[All initial-length tails and their signs]
\label{thm:tail}
For $n\ge3$ and $1\le d\le n$,
\begin{equation}
\rho(n)\equiv\delta_nE_d\pmod{10^d},\qquad
\delta_n=\begin{cases}-1,&n\equiv1\pmod4,\\+1,&n\not\equiv1\pmod4.
\end{cases}
\label{eq:tail}
\end{equation}
In particular, for every $n\ge10$ the last ten digits of $\rho(n)$ are
\begin{equation}
\begin{cases}
1787109375,&n\equiv1\pmod4,\\
8212890625,&n\not\equiv1\pmod4.
\end{cases}
\label{eq:ten}
\end{equation}
\end{theorem}
\begin{proof}
Because $5\mid r(n)$, $5^n\mid\rho(n)$.
If $n$ is even, apply \eqref{eq:lte2general} to $r(n)^n$:
\[
\nu_2(\rho(n)-1)=c_2(r(n))+\nu_2(n)=s+\nu_2(n)=n.
\]
Hence $\rho(n)\equiv1\pmod{2^n}$.
If $n$ is odd, $s=n$ and \eqref{eq:main} becomes
\[
r(n)=\begin{cases}3\cdot2^n-1,&n\equiv1\pmod4,\\
3\cdot2^n+1,&n\equiv3\pmod4.
\end{cases}
\]
Raising to the odd exponent $n$ preserves the indicated sign modulo
$2^n$. Thus in all cases $\rho(n)\equiv\delta_n\pmod{2^n}$.
Apply the Chinese remainder theorem and reduce to any $d\le n$.
Finally, direct calculation from \eqref{eq:Edformula} gives
$E_{10}=8212890625$ and $10^{10}-E_{10}=1787109375$.
\end{proof}

The ten-digit assertion is stronger than simply choosing one of two
possible suffixes: it determines which suffix occurs from $n\bmod4$.
It also starts at $n=10$, one index earlier than the range $n>10$
specified in the question \cite{mo2025}.

There are two distinct limiting notions in this discussion. The
compatible residues $E_d$ describe a $10$-adic integer with components
$(1,0)$ in $\Z_2\times\Z_5$. They are not a claim that the ordinary real
sequence $\rho(n)$ converges. They describe the minimizing \emph{bases},
not the eventual tail of a tower formed from one fixed base.

\section{A two-parameter extremal theorem}
\label{sec:twoparameter}

The proof is not confined to the diagonal where degree and speed coincide.
For positive integers $k,m$, put
\[
R_k(m)=\min\{x>1:10\nmid x,\ V(x^k)=m\}
\]
whenever this set is nonempty. Here $k$ is the power degree and $m$ is
the requested speed.

\begin{theorem}[An explicit region of global minimality]
\label{thm:two}
Let
\[
s=m-\nu_2(k)\ge2,\qquad t=m-\nu_5(k)\ge1.
\]
If
\begin{equation}
5^t>B_s^2+1,
\label{eq:twocriterion}
\end{equation}
then $R_k(m)=B_s$, and the least $k$th perfect power of speed $m$ is
$B_s^k$.
\end{theorem}
\begin{proof}
By \eqref{eq:powertransfer}, $B_s$ has exactly the requested speed after
raising to the $k$th power. It is least among roots divisible by $5$.
Any competing root not divisible by $5$ has $c_5(x)\ge t$, so
\cref{lem:barrier} and \eqref{eq:twocriterion} rule it out.
\end{proof}

\begin{corollary}[Fixed-degree eventual formula]
\label[corollary]{cor:fixdegree}
For every fixed $k$, the equality
\[
R_k(m)=B_{m-\nu_2(k)}
\]
holds for all sufficiently large $m$. A directly checkable sufficient
condition, in addition to $s\ge2,t\ge1$, is
\begin{equation}
\left(\frac54\right)^m\ge
16\,\frac{5^{\nu_5(k)}}{4^{\nu_2(k)}}.
\label{eq:explicitwedge}
\end{equation}
\end{corollary}
\begin{proof}
For $s\ge2$,
\[
B_s^2+1\le(3\cdot2^s+1)^2+1<16\cdot4^s.
\]
Condition \eqref{eq:explicitwedge} gives
$5^{m-\nu_5(k)}\ge16\cdot4^{m-\nu_2(k)}>B_s^2+1$.
For fixed $k$, the left side of \eqref{eq:explicitwedge} tends to
infinity, while the right side is constant.
\end{proof}

This result gives an effective answer in an infinite region of the
$(k,m)$ plane. It does not assert that the displayed inequality is
necessary, nor does it claim a simple formula for every small pair
outside that region.

\section{Exact densities and what the heuristic was seeing}
\label{sec:density}

A density argument is not a minimality proof: a sparse set can still
contain the smallest element. Nevertheless, the local formulas make
precise the statistical preference for roots divisible by $5$.
For a subset $A\subseteq\N$, its natural density, when it exists, is
\[
\lim_{X\to\infty}\frac{\#(A\cap[1,X])}{X}.
\]
All densities below exist because the conditions, apart from the single
excluded integer $1$, are finite unions of residue classes.

\begin{theorem}[Tail and exact-speed densities]
\label{thm:density}
Fix $k,m\ge1$ with $s=m-\nu_2(k)\ge2$ and $t=m-\nu_5(k)\ge1$.
The density, among all positive integers $x$, of
\[
\{x>1:10\nmid x,\ V(x^k)\ge m\}
\]
is
\begin{equation}
\mathcal D_k(m)=\frac{2^{1-s}}5+\frac2{5^t}
+\frac{2^{3-s}}{5^t}.
\label{eq:densitytail}
\end{equation}
The density of the corresponding set with $V(x^k)=m$ is
\begin{equation}
\mathcal E_k(m)=\frac{2^{-s}}5+\frac8{5^{t+1}}
+\frac{9\cdot2^{2-s}}{5^{t+1}}.
\label{eq:densityexact}
\end{equation}
\end{theorem}
\begin{proof}
For odd $x$, $c_2(x)\ge s$ is equivalent to
$x\equiv\pm1\pmod{2^s}$, which occupies $2/2^s$ of all integers.
For $5\nmid x$, $c_5(x)\ge t$ is equivalent to $x^4\equiv1\pmod{5^t}$.
There are exactly four roots: $1,-1$, and the two roots of $-1$ from
\cref{lem:hensel}.

Partition the permitted integers into three disjoint types. Odd
multiples of $5$ contribute
$(2/2^s)(1/5)=2^{1-s}/5$. Even integers not divisible by $5$ contribute
$(1/2)(4/5^t)=2/5^t$. Integers coprime to $10$ contribute
$(2/2^s)(4/5^t)=2^{3-s}/5^t$. The products follow from the Chinese
remainder theorem. Their sum proves \eqref{eq:densitytail}.
For exact speed, subtract the density at threshold $m+1$, which replaces
$s,t$ by $s+1,t+1$. This gives \eqref{eq:densityexact}.
\end{proof}

These are densities among \emph{roots} $x$, not among the much sparser
set of $k$th powers as numbers on the ordinary number line. Also the
denominator includes all positive integers; conditioning on $10\nmid x$
would divide the displayed densities by $9/10$.

For the diagonal $k=m=n$, the contribution $2^{-s}/5$ to exact-speed
density comes from odd multiples of $5$. The ratio of the other two
contributions to this one is exactly
\[
\frac{8\cdot2^s+36}{5^t}.
\]
Since $2^s\le2^n$ and $5^t\ge5^n/n$, this ratio tends to zero.
Thus almost all eligible roots, in the density sense as the parameter
$n$ grows, belong to the final-digit-$5$ class. This explains the
heuristic but is logically separate from the proof that the
\emph{least} root always belongs to that class for $n\ge3$.

\section{Asymptotics and oscillations of the minimum}
\label{sec:asymptotics}

Write
\[
c_s=\begin{cases}1,&s\text{ even},\\3,&s\text{ odd},\end{cases}
\qquad
\varepsilon_s=\begin{cases}-1,&s\equiv0,1\pmod4,\\+1,&s\equiv2,3\pmod4.
\end{cases}
\]
Then $r(n)=c_s2^s+\varepsilon_s$ for $n\ge3$.

\begin{proposition}[Logarithmic asymptotics]
\label[proposition]{prop:asymp}
As $n\to\infty$,
\begin{equation}
\log\rho(n)
=n\bigl(n-\nu_2(n)\bigr)\log2+n\log c_s
+O\bigl(n2^{-s}\bigr).
\label{eq:logasymptotic}
\end{equation}
Consequently
\begin{equation}
\lim_{n\to\infty}\rho(n)^{1/n^2}=2,
\qquad
\log_2\rho(n)=n^2+O(n\log n).
\label{eq:rootlimit}
\end{equation}
\end{proposition}
\begin{proof}
Factor $r(n)=c_s2^s(1+\varepsilon_s/(c_s2^s))$ and use
$\log(1+u)=O(u)$ uniformly for $|u|\le1/4$. Multiply by $n$.
Since $\nu_2(n)\le\log_2 n$, we have
$s\ge n-\log_2n\to\infty$. Dividing \eqref{eq:logasymptotic} by $n^2$
then proves both conclusions.
\end{proof}

Despite the clean quadratic logarithmic growth, the root sequence is
not monotone. Some adjacent values coincide; others decrease sharply.
For example, if $n=4j+1\ge5$, then
\[
n-\nu_2(n)=n=(n+1)-\nu_2(n+1),
\]
so $r(4j+1)=r(4j+2)$.

\begin{proposition}[Complete accumulation set of normalized roots]
\label[proposition]{prop:accumulation}
The set of subsequential limits of $r(n)/2^n$ is
\begin{equation}
\{0,3\}\ \cup\
\{3/2^j:j\ge1\text{ odd}\}\ \cup\
\{1/2^j:j\ge2\text{ even}\}.
\label{eq:accumulation}
\end{equation}
In particular, $\liminf r(n)/2^n=0$ and $\limsup r(n)/2^n=3$.
\end{proposition}
\begin{proof}
Along odd $n$, $s=n$ is odd and the normalized root tends to $3$.
Along integers with fixed $j=\nu_2(n)\ge1$, the integer $n$ is even, so
$s=n-j$ is odd precisely when $j$ is odd. Hence
\[
\frac{r(n)}{2^n}=\frac{c_s}{2^j}+O(2^{-n})
\]
has the indicated limit. Each fixed valuation class contains infinitely
many integers. Along $n=2^j$ with $j\to\infty$, the ratio tends to zero.
Conversely, along a subsequence of indices tending to infinity, either
the $2$-adic valuations have an unbounded subsequence, forcing a further
limit of zero, or some fixed valuation occurs infinitely often. If the
original subsequence converges, its limit must be one of these values.
\end{proof}

The limit $\rho(n)^{1/n^2}\to2$ concerns a sequence of specially chosen
\emph{bases}. It should not be confused with a bound on the growth of
$t\mapsto a\uparrow\uparrow t$, which is a different variable and a
very different sequence.

\section{Algorithms, validation, and reproducibility}
\label{sec:computing}

\subsection{The constant-size formula is not a tower computation}
Once $n$ is given, the formula requires the valuation $\nu_2(n)$, a
residue modulo $4$, a binary shift, and one small multiplication and
addition. The full root $r(n)$ has $\Theta(n)$ bits, while $\rho(n)$ has
$\Theta(n^2)$ bits. Thus writing the root or the power has an unavoidable
output cost. Describing the formula as ``constant time'' would be
misleading when the input $n$ is encoded in binary.

A direct implementation of the decisive part is:
\begin{lstlisting}[language=Python]
def minimum_root(n: int) -> int:
    if n < 1:
        raise ValueError("n must be positive")
    if n == 1:
        return 2
    if n == 2:
        return 7
    v2 = (n & -n).bit_length() - 1
    s = n - v2
    coefficient = 3 if s % 2 else 1
    sign = -1 if s % 4 in (0, 1) else 1
    return coefficient * (1 << s) + sign
\end{lstlisting}
The accompanying module includes argument validation, speed evaluation,
modular tails, the exact-degree variant, and command-line support.
It uses Python's standard library only. To compute the full answer,
one evaluates \texttt{pow(minimum\_root(n), n)}. Repeated squaring uses
$O(\log n)$ integer multiplications; their sizes must be included in a
bit-complexity estimate.

For only the last $d\le n$ decimal digits with $n\ge3$, use
\cref{thm:tail}; it does not require materializing a root of $\Theta(n)$
bits. Formula \eqref{eq:Edformula} involves integers with $O(d)$ bits and
one modular inverse, followed by the sign test on $n\bmod4$.

\subsection{What was actually verified}
The included run was executed with Python 3.13.5. All tests passed.
The detailed machine-readable record is \texttt{data/verification.json}.
The checks deliberately use more than one computational route.

First, the compact speed formula was compared with a separate
transcription of the published residue-class formula for every
$a=2,\ldots,20000$ not divisible by $10$: 17,999 comparisons.
Second, the power-transfer formula was checked against actual integer
powers and that separate formula for $2\le x\le500$, $10\nmid x$,
and $1\le k\le30$: 13,470 comparisons.

Third, unrestricted minima were exhaustively searched for
$1\le n\le20$, using actual values $x^n$ rather than merely checking
the proposed root. Exact-degree minima were exhaustively searched for
$1\le n\le15$. Every proposed root was independently checked to have
speed $n$ through $n=10000$. The first 50 roots agree with the values
reported by Alekseyev in the source discussion.

Fourth, the integer comparisons in \cref{tab:small} and the Hensel
certificates in \cref{tab:certs} were regenerated exactly. Modular tail
identities were checked for $3\le n\le2000$ at lengths
$1,\min(n,10),\min(n,30)$. A separate perfect-power test checked the
assertion about $B_s$ for $2\le s\le200$.

Finally, tower residues were computed by an independent modular
exponentiation routine for 17 representative bases, through tower
height $14$, with 160 decimal digits of modulus. The resulting
$D_0,\ldots,D_{13}$ agree with the predicted eventual speeds. The tower
routine itself was first checked in 88 comparisons with literal towers
small enough to construct. Several complete residue periods were also
enumerated to validate the density formula.

The role of these computations is verification and error detection.
The infinite quantifier in \cref{thm:main} is discharged by
\cref{lem:cutoff}, not by checking roots through $n=10000$.

\subsection{Reproduction commands}
From the archive's top-level directory, run
\begin{lstlisting}[language=bash]
python code/verify.py
python code/tetration_minima.py 15 --root-only
python code/tetration_minima.py 3 --strict
python code/tetration_minima.py 100 --digits 10
python code/tetration_minima.py --table roots.csv --count 1000
pdflatex -interaction=nonstopmode -halt-on-error article.tex
pdflatex -interaction=nonstopmode -halt-on-error article.tex
\end{lstlisting}
The second command prints $98305$; the third prints $166375$.
The fourth prints \texttt{8212890625}. Running the verification script
regenerates the CSV files and verification record, including the
finite certificates. No network access, symbolic algebra system, or
floating-point computation is required by the mathematical checks.
Elapsed time in the log is descriptive of that particular run, not a
portable performance guarantee.

\section{What has been settled and what has not}
\label{sec:status}

For the definitions in \cref{sec:problem}, the global minimum is now
given by \eqref{eq:main} for every positive $n$. The proof also determines
the maximal-degree variant and the precise ten-digit-tail choice.
The unrestricted minimum and its exact-degree strengthening are not
conditional on a conjecture or an unverified random-digit assumption.

The historical claim is narrower: the specific conjectural pattern in
the consulted MathOverflow question and answer is proved here. The
underlying speed formula, the method of using valuations, the observed
initial values, and the idea that the final-digit-$5$ branch should
win all have prior provenance. This report's contribution is the
explicit global argument based on the quadratic barrier, together with
the closed formula and the derived refinements. The sources checked
are not an exhaustive review of all unpublished notes or discussions.
Independent review remains appropriate before making a formal priority
claim or submitting the result for publication.

The two-parameter theorem leaves small pairs outside its sufficient
region untreated, and no arbitrary-radix analogue is proved. A natural
next problem is to determine which radices admit a similarly simple
extremal branch. The crucial feature here is that a $5$-adic fourth-root
condition factors into pieces of degree at most two, making a
$5^{m/2}$ size barrier available. In another radix, the relevant
cyclotomic factors and competing prime-adic slopes can change the
comparison. Such an extension would require its own proof rather than
a formal replacement of the number $10$.

\appendix
\section{A short proof suitable for checking independently}
\label{app:short}

For a reader willing to use the known congruence-speed formula, the
entire main argument compresses to the following.
For odd $x>1$, put $c_2(x)=\nu_2(x^2-1)-1$; for $5\nmid x>1$, put
$c_5(x)=\nu_5(x^4-1)$. Lifting gives
$c_2(x^n)=c_2(x)+\nu_2(n)$ and
$c_5(x^n)=c_5(x)+\nu_5(n)$.

Set $s=n-\nu_2(n)$ and $t=n-\nu_5(n)$. The least odd multiple of $5$
with $c_2(x)=s\ge2$ is $B_s$, by solving
$x=j2^s\pm1$ with $j$ odd and then reducing modulo $5$.
It yields $V(B_s^n)=n$.

Any other competing root must satisfy $c_5(x)\ge t$. Since exactly one
of $x-1,x+1,x^2+1$ is divisible by $5$, it follows that
$x^2+1\ge5^t$. Hence $5^t>B_s^2+1$ proves global minimality.
For $n\ge25$, use
\[
5^t\ge5^n/n>(3\cdot2^n+1)^2+1\ge B_s^2+1;
\]
the middle inequality follows from its checked value at $25$ and the
ratio argument in \cref{lem:cutoff}.

For $3\le n\le24$, the same inequality fails only at
$3,5,7,9,10,15$. At those indices, use respectively the smaller roots
$57,182,1068,32318,32318,110443$ of $-1$ modulo
$5^3,5^4,5^5,5^7,5^7,5^8$. They and the least nontrivial $\pm1$
residues exceed $B_s$. This excludes every remaining competitor.
Checking $n=1,2$ gives $2$ and $7$.

This summary uses only finitely many small, explicit arithmetic
certificates. The longer article proves the speed formula and the
lifting steps as well.

\section{Tables and archive inventory}
\label{app:data}

\subsection{First twenty-five minimizing roots}
\begin{center}
\begin{tabular}{rr@{\qquad\qquad}rr@{\qquad\qquad}rr}
\toprule
$n$&$r(n)$&$n$&$r(n)$&$n$&$r(n)$\\
\midrule
1&2&10&1535&19&1572865\\
2&7&11&6145&20&262145\\
3&25&12&1025&21&6291455\\
4&5&13&24575&22&6291455\\
5&95&14&24575&23&25165825\\
6&95&15&98305&24&6291455\\
7&385&16&4095&25&100663295\\
8&95&17&393215&&\\
9&1535&18&393215&&\\
\bottomrule
\end{tabular}
\end{center}
For the exact-degree sequence, replace only the root at $n=3$ by $55$.
The desired base is the $n$th power of the root in this table.

\subsection{A lifting chain for square roots of \texorpdfstring{$-1$}{-1}}
The following chain starts at $2\pmod5$ and follows its unique lift.
The other root is the negative modulo $5^k$.
\begin{center}
\begin{tabular}{rrrr}
\toprule
$k$&$5^k$&lift of $2\pmod5$&smaller root $d_k$\\
\midrule
1&5&2&2\\
2&25&7&7\\
3&125&57&57\\
4&625&182&182\\
5&3125&2057&1068\\
6&15625&14557&1068\\
7&78125&45807&32318\\
8&390625&280182&110443\\
\bottomrule
\end{tabular}
\end{center}

\subsection{Included files}
\begin{description}[style=nextline,leftmargin=1em,labelwidth=0pt,labelindent=0pt]
\item[\texttt{article.tex}, \texttt{article.pdf}]
The complete source and rendered article.
\item[\texttt{code/tetration\_minima.py}]
The exact formulas, speed functions, Hensel lifts, tails, strict variant,
and command-line interface.
\item[\texttt{code/verify.py}]
Independent exhaustive checks and generation of all finite certificates.
\item[\texttt{data/minimum\_roots\_1\_1000.csv}]
The first 1000 roots, their valuation parameters, and checked speeds.
\item[\texttt{data/rho\_1\_50.txt}]
The first 50 complete integers $\rho(n)$ in index--value format.
\item[\texttt{data/finite\_certificates.csv}, \texttt{data/small\_range.csv}]
Machine-readable versions of the finite proof calculations.
\item[\texttt{data/verification.json}, \texttt{data/verification.log}]
The validation results, including independent tower checks.
\item[\texttt{README.md}, \texttt{sources.md}, \texttt{build.sh}]
Usage notes, source provenance and status, and a simple build command.
\end{description}

\begin{thebibliography}{9}
\bibitem{ripa2020}
M.~Rip\`a,
\emph{On the constant congruence speed of tetration},
Notes on Number Theory and Discrete Mathematics \textbf{26} (2020),
no.~3, 245--260.
\href{https://doi.org/10.7546/nntdm.2020.26.3.245-260}
{doi:10.7546/nntdm.2020.26.3.245-260}.

\bibitem{ripa2021}
M.~Rip\`a,
\emph{The congruence speed formula},
Notes on Number Theory and Discrete Mathematics \textbf{27} (2021),
no.~4, 43--61.
\href{https://doi.org/10.7546/nntdm.2021.27.4.43-61}
{doi:10.7546/nntdm.2021.27.4.43-61}.
Also \href{https://arxiv.org/abs/2208.02622}{arXiv:2208.02622} (2022).

\bibitem{ripaonnis2022}
M.~Rip\`a and L.~Onnis,
\emph{Number of stable digits of any integer tetration},
Notes on Number Theory and Discrete Mathematics \textbf{28} (2022),
no.~3, 441--457.
\href{https://doi.org/10.7546/nntdm.2022.28.3.441-457}
{doi:10.7546/nntdm.2022.28.3.441-457}.

\bibitem{mo2025}
M.~Rip\`a (question) and M.~Alekseyev (answer),
\emph{Closed form for the general term of
$2,49,15625,625,\ldots$},
MathOverflow, question 487698, 12--14 February 2025.
\url{https://mathoverflow.net/questions/487698/}.
Question, answer, and associated conjectural statements checked
19 September 2026.

\bibitem{ripa2026}
M.~Rip\`a,
\emph{On the relation between perfect powers and tetration frozen digits},
Journal of AppliedMath \textbf{2} (2024), no.~5, Article 1771.
\href{https://doi.org/10.59400/jam1771}{doi:10.59400/jam1771}.
Also \href{https://arxiv.org/abs/2602.00252}{arXiv:2602.00252v1},
30 January 2026.
Cited for the existence program and context, not as a proof of the
extremal theorem established here.
\end{thebibliography}
\end{document}
